\section{Resultados da revisão}
\label{sec:resultados}

A Tabela~\ref{tab:estudos} identifica os \nIncluidos{} estudos primários
analisados em texto integral (Secção~\ref{sec:selecao}), numerados por ordem
cronológica (E01--E19), com o tipo de publicação e a categoria que lhes foi
atribuída na taxonomia da Secção~\ref{sec:taxonomia}.

\begin{table}[ht]
  \centering
  \caption{Estudos primários incluídos. Tipo: R revista, C conferência, W
    \emph{workshop}. Categorias: ver Figura~\ref{fig:taxonomia}.}
  \label{tab:estudos}
  \footnotesize
  \setlength{\tabcolsep}{4pt}
  \begin{tabular}{@{}llccc@{\hspace{1.2em}}llccc@{}}
    \toprule
    \textbf{ID} & \textbf{Estudo} & \textbf{Ano} & \textbf{Tipo} & \textbf{Cat.} &
    \textbf{ID} & \textbf{Estudo} & \textbf{Ano} & \textbf{Tipo} & \textbf{Cat.} \\
    \midrule
    \E{01} & Singh et al.~\cite{singh2015secure}               & 2015 & C & C4 &
    \E{11} & Vaccari et al.~\cite{vaccari2020mqttset}          & 2020 & R & C7 \\
    \E{02} & Niruntasukrat et al.~\cite{niruntasukrat2016authorization} & 2016 & W & C3 &
    \E{12} & Segarra et al.~\cite{segarra2020mqttz}            & 2020 & C & C6 \\
    \E{03} & Rizzardi et al.~\cite{rizzardi2016aups}           & 2016 & R & C3 &
    \E{13} & Dahlmanns et al.~\cite{dahlmanns2021entrust}      & 2021 & W & C5 \\
    \E{04} & Andy et al.~\cite{andy2017attack}                 & 2017 & C & C1 &
    \E{14} & Spielvogel et al.~\cite{spielvogel2021tls}        & 2021 & W & C6 \\
    \E{05} & Calabretta et al.~\cite{calabretta2018mqttauth}   & 2018 & R & C3 &
    \E{15} & Prantl et al.~\cite{prantl2021performance}        & 2021 & C & C6 \\
    \E{06} & Hernández Ramos et al.~\cite{hernandez2018fuzzing} & 2018 & R & C1 &
    \E{16} & Van Glabbeek et al.~\cite{vanglabbeek2022flexible} & 2022 & R & C5 \\
    \E{07} & Dinculeană e Cheng~\cite{dinculeana2019vulnerabilities} & 2019 & R & C5 &
    \E{17} & Buccafurri et al.~\cite{buccafurri2023mqtta}      & 2023 & R & C8 \\
    \E{08} & Malina et al.~\cite{malina2019secure}             & 2019 & C & C5 &
    \E{18} & Alotaibi et al.~\cite{alotaibi2024secure}         & 2024 & R & C7 \\
    \E{09} & Sadio et al.~\cite{sadio2019lightweight}          & 2019 & C & C5 &
    \E{19} & Javadi et al.~\cite{javadi2026cmap}               & 2026 & R & C2 \\
    \E{10} & Vaccari et al.~\cite{vaccari2020slowite}          & 2020 & R & C1 & & & & & \\
    \bottomrule
  \end{tabular}
\end{table}

% ---------------------------------------------------------------------------
\subsection{Caracterização dos estudos (QE1, QE2)}
\label{sec:estatistica}

\paragraph{Publicações por ano (QE1).} A Figura~\ref{fig:por-ano} mostra a
distribuição temporal dos \nElegiveis{} estudos elegíveis após a triagem. O
interesse surge logo após a normalização do MQTT~3.1.1 (2014) e cresce de forma
contínua: os anos de 2023 a 2026 reúnem 410 estudos ($56{,}9\,\%$), e o ano de 2026
está contado só até à data das pesquisas. Destes estudos elegíveis, só
\nIncluidos{} foram analisados em texto integral (Secção~\ref{sec:selecao}); as
restantes estatísticas desta secção referem-se a esse conjunto.

\begin{figure}[ht]
  \centering
  \begin{tikzpicture}
    \begin{axis}[
        ybar, bar width=12pt, width=0.92\textwidth, height=5.2cm,
        ymin=0, ymax=130, ytick={0,20,...,120},
        symbolic x coords={2014,2015,2016,2017,2018,2019,2020,2021,2022,2023,2024,2025,2026},
        xtick=data, x tick label style={font=\small},
        ylabel={N.º de estudos}, ylabel style={font=\small},
        nodes near coords, nodes near coords style={font=\scriptsize},
        axis lines*=left, ymajorgrids, grid style={cinza!25},
        enlarge x limits=0.05]
      \addplot[fill=acento!70, draw=iscte] coordinates {
        (2014,2) (2015,4) (2016,9) (2017,17) (2018,30) (2019,58) (2020,51)
        (2021,74) (2022,66) (2023,75) (2024,115) (2025,101) (2026,119)};
    \end{axis}
  \end{tikzpicture}
  \caption{Estudos elegíveis após a triagem, por ano de publicação ($n = \nElegiveis$; 2026 até à data das pesquisas). Destes, \nIncluidos{} foram analisados em texto integral.}
  \label{fig:por-ano}
\end{figure}

\paragraph{Tipo de publicação (QE2).} Dos \nIncluidos{} estudos analisados, dez ($52{,}6\,\%$) foram publicados
em revistas, seis ($31{,}6\,\%$) em conferências e três ($15{,}8\,\%$) em
\emph{workshops}.

% ---------------------------------------------------------------------------
\subsection{Taxonomia das soluções (QF2)}
\label{sec:taxonomia}

Os estudos foram agrupados em oito categorias segundo o ponto de atuação e o
objetivo do seu contributo principal (Figura~\ref{fig:taxonomia}).

\begin{figure}[ht]
  \centering
  \begin{tikzpicture}[
      raiz/.style={draw=iscte, fill=iscte, text=white, rounded corners=2pt,
                   font=\small\bfseries, minimum height=0.8cm, text width=4cm, align=center},
      ramo/.style={draw=iscte, fill=fundotab, rounded corners=2pt, font=\scriptsize\bfseries,
                   text width=3.3cm, minimum height=0.7cm, align=center},
      folha/.style={draw=cinza, rounded corners=2pt, font=\scriptsize, text width=3.3cm,
                    align=left, minimum height=0.6cm},
      ligacao/.style={draw=iscte, thick}]
    \node[raiz] (r) {Segurança em MQTT ($n=\nIncluidos$)};
    \node[ramo, below left=1.0cm and 2.0cm of r] (a) {Compreender ameaças};
    \node[ramo, below=1.0cm of r] (b) {Proteger ligação e \emph{broker}};
    \node[ramo, below right=1.0cm and 2.0cm of r] (c) {Proteger a mensagem e os dados};
    \node[folha, below=0.35cm of a] (c1) {C1 Análise de vulnerabilidades e ataques (3)};
    \node[folha, below=0.2cm of c1] (c7) {C7 Deteção de intrusões (2)};
    \node[folha, below=0.35cm of b] (c2) {C2 Autenticação (1)};
    \node[folha, below=0.2cm of c2] (c3) {C3 Autorização e controlo de acesso (3)};
    \node[folha, below=0.2cm of c3] (c6) {C6 Proteção do \emph{broker} e do transporte (3)};
    \node[folha, below=0.35cm of c] (c5) {C5 Proteção leve fim-a-fim da mensagem (5)};
    \node[folha, below=0.2cm of c5] (c4) {C4 Cifra baseada em atributos (1)};
    \node[folha, below=0.2cm of c4] (c8) {C8 Privacidade e anonimato (1)};
    \draw[ligacao] (r.south) -- ++(0,-0.4) -| (a.north);
    \draw[ligacao] (r.south) -- (b.north);
    \draw[ligacao] (r.south) -- ++(0,-0.4) -| (c.north);
    \foreach \p/\f in {a/c1, b/c2, c/c5} \draw[ligacao] (\p) -- (\f);
    \draw[ligacao] (c1) -- (c7); \draw[ligacao] (c2) -- (c3); \draw[ligacao] (c3) -- (c6);
    \draw[ligacao] (c5) -- (c4); \draw[ligacao] (c4) -- (c8);
  \end{tikzpicture}
  \caption{Taxonomia dos estudos incluídos, com o número de estudos por categoria.}
  \label{fig:taxonomia}
\end{figure}

\subsubsection{C1 --- Análise de vulnerabilidades e ataques}

Andy et al.~\cite{andy2017attack} mostram que, por omissão, o MQTT não cifra os
dados e demonstram ataques à privacidade, à autenticação e à integridade. Hernández
Ramos et al.~\cite{hernandez2018fuzzing} propõem uma técnica de \emph{fuzzing}
baseada em \emph{templates} que encontrou falhas novas em implementações reais. Vaccari et
al.~\cite{vaccari2020slowite} propõem o SlowITe, um ataque de negação de serviço de
baixa taxa, eficaz contra serviços reais com e sem cifra.

\subsubsection{C2 --- Autenticação}

Javadi et al.~\cite{javadi2026cmap} propõem o CMAP, autenticação mútua sem
certificados, sobre curvas elípticas, entre cada dispositivo e o \emph{broker}. A
segurança é demonstrada no modelo do oráculo aleatório e o protocolo é avaliado em
Raspberry Pi (338~ms de latência de autenticação contra 713~ms com TLS) e em
simulação NS-3. A resistência à repetição que reportam diz respeito às mensagens de
autenticação, protegidas por \emph{nonces} frescos, e não às mensagens publicadas.

\subsubsection{C3 --- Autorização e controlo de acesso}

Niruntasukrat et al.~\cite{niruntasukrat2016authorization} adaptam o OAuth~1.0a a
nós restritos e mostram atrasos de autorização mínimos numa plataforma real.
Rizzardi et al.~\cite{rizzardi2016aups} estendem o MQTT com gestão de chaves e
aplicação de políticas (AUPS, em código aberto). Calabretta et
al.~\cite{calabretta2018mqttauth} propõem o MQTT-Auth: o AugPAKE garante a
confidencialidade e dois \emph{tokens} autenticam e autorizam o acesso a cada
tópico.

\subsubsection{C4 --- Cifra baseada em atributos}

Singh et al.~\cite{singh2015secure} propõem o SMQTT e o SMQTT-SN, que acrescentam
ao MQTT cifra baseada em atributos com política na chave ou no criptograma
(KP/CP-ABE) sobre criptografia de curvas elípticas (ECC) leve, avaliada em simulação: só os subscritores com
atributos compatíveis decifram. A ABE resolve a confidencialidade de grupo, mas com
um custo computacional elevado para dispositivos restritos.

\subsubsection{C5 --- Proteção leve fim-a-fim da mensagem}

É a categoria mais numerosa. Dinculeană e Cheng~\cite{dinculeana2019vulnerabilities}
enviam valores sob a forma de HMAC (\emph{Value-to-HMAC}), com melhor desempenho do
que uma cifra simétrica; os próprios autores remetem para outros mecanismos a
autenticação dos clientes, a distribuição de chaves e a proteção contra
repetição~\cite[§6]{dinculeana2019vulnerabilities}. Malina et
al.~\cite{malina2019secure} propõem três níveis de segurança, do simples controlo de
adulteração à autenticação mútua de todas as partes. Todos assinam com Schnorr uma
entrada que inclui o nível, o tópico, a identidade do publicador, os dados e um
carimbo temporal~\cite[§4.5, Eq.~6]{malina2019secure}, o que dá não repúdio e liga a
mensagem ao tópico. Nos níveis 2 e 3 o carimbo serve de proteção contra
repetição~\cite[§5.2]{malina2019secure}, mas o \emph{broker} decifra o
conteúdo~\cite[§4.6]{malina2019secure}. Sadio et al.~\cite{sadio2019lightweight} usam
ChaCha20-Poly1305 fim-a-fim, com um protótipo em Arduino UNO. Dahlmanns et
al.~\cite{dahlmanns2021entrust} propõem o ENTRUST: um servidor de chaves distribui,
segundo uma ACL, uma chave por tópico aos clientes autorizados, sem alterar o
protocolo nem confiar no \emph{broker}~\cite[§5.2--5.3]{dahlmanns2021entrust}. Cada
mensagem leva um contador por publicador, que permite detetar duplicados e
perdas~\cite[§5.3]{dahlmanns2021entrust}. Van Glabbeek et
al.~\cite{vanglabbeek2022flexible} propõem três opções (autenticação e não repúdio;
mais confidencialidade; mais anonimato), todas com assinatura Schnorr e carimbo
temporal, com um \emph{broker} honesto mas curioso e uma autoridade de
certificação (CA) totalmente confiável~\cite[§3.1, §4.3]{vanglabbeek2022flexible}. Face ao MQTT sobre TLS, reduzem
em 80\,\% o \emph{overhead} de comunicação e em 40\,\% o computacional no
estabelecimento da ligação.

\subsubsection{C6 --- Proteção do \emph{broker} e do transporte}

Segarra et al.~\cite{segarra2020mqttz} executam o \emph{broker} num ambiente de
execução confiável (ARM TrustZone), avaliado com cargas reais de uma aplicação
médica. Spielvogel et al.~\cite{spielvogel2021tls} observam que um \emph{broker}
reencaminha mensagens independentemente da proteção de cada salto e propõem que
publicadores e subscritores possam exigir um nível de segurança mínimo. Prantl et
al.~\cite{prantl2021performance} medem o impacto do TLS num microcontrolador e
concluem que depende sobretudo das condições da rede e da frequência de
restabelecimento de ligações.

\subsubsection{C7 --- Deteção de intrusões}

Vaccari et al.~\cite{vaccari2020mqttset} publicam o MQTTset, um conjunto de dados
de tráfego MQTT legítimo e de ataques para treinar detetores. Alotaibi et
al.~\cite{alotaibi2024secure} detetam anomalias em tempo real com um enquadramento
distribuído (H2O), com os melhores resultados obtidos por XGBoost.

\subsubsection{C8 --- Privacidade e anonimato}

Buccafurri et al.~\cite{buccafurri2023mqtta} propõem o MQTT-A, uma rede de
\emph{brokers} intermédios que impede o \emph{broker} final de identificar a origem
e o destino das mensagens, sem alterar a infraestrutura MQTT.

% ---------------------------------------------------------------------------
\subsection{Comparação dos mecanismos (QF3, QF4)}
\label{sec:comparacao}

A Tabela~\ref{tab:comparacao} compara os 16 estudos que propõem mecanismos
(C2 a C8); os três estudos de análise (C1) ficam de fora. Cada célula foi
classificada a partir do texto integral, e a secção ou página que a sustenta está
registada no ficheiro de rastreabilidade do repositório~\cite{mqttsecure2026repo}.
\enquote{Não reportado} (\Nr) não equivale a \enquote{ausente por desenho}
(\Nao).

\begin{table}[p]
  \centering
  \caption{Comparação dos estudos que propõem mecanismos. Propriedades: C
    confidencialidade, I integridade, Au autenticação, Az autorização, V validade
    temporal, R anti-repetição, NR não repúdio, D disponibilidade/deteção, P privacidade.
    BnC: funciona com \emph{broker} não confiável. Aval.: E experimental, F formal,
    S simulação, A analítica. Restr.: avaliado ou desenhado para dispositivos
    restritos. \Sim{} garantido; \Parc{} parcialmente; \Nao{} não, por desenho;
    \Nr{} não reportado.}
  \label{tab:comparacao}
  \scriptsize
  \setlength{\tabcolsep}{2.6pt}
  \renewcommand{\arraystretch}{1.25}
  \begin{tabular}{@{}lL{2.55cm}L{3.55cm}ccccccccc|ccc@{}}
    \toprule
    \textbf{ID} & \textbf{Estudo} & \textbf{Mecanismo principal}
      & \textbf{C} & \textbf{I} & \textbf{Au} & \textbf{Az} & \textbf{V} & \textbf{R} & \textbf{NR}
      & \textbf{D} & \textbf{P} & \textbf{BnC} & \textbf{Aval.} & \textbf{Restr.} \\
    \midrule
    \multicolumn{15}{@{}l}{\textit{C2 --- Autenticação}} \\
    \E{19} & Javadi et al.~\cite{javadi2026cmap}            & Autenticação mútua sem certificados (ECC)
      & \Parc & \Parc & \Sim & \Nr & \Nr & \Nr & \Nr & \Nr & \Nr & \Nao & F, E, S & \Sim \\
    \multicolumn{15}{@{}l}{\textit{C3 --- Autorização e controlo de acesso}} \\
    \E{02} & Niruntasukrat et al.~\cite{niruntasukrat2016authorization} & OAuth 1.0a adaptado
      & \Nr & \Nr & \Parc & \Sim & \Nr & \Nr & \Nr & \Nr & \Nr & \Nr & E & \Sim \\
    \E{03} & Rizzardi et al.~\cite{rizzardi2016aups}        & Gestão de chaves e políticas (AUPS)
      & \Parc & \Nr & \Nr & \Sim & \Nr & \Nr & \Nr & \Nr & \Nr & \Nr & A, E & \Sim \\
    \E{05} & Calabretta et al.~\cite{calabretta2018mqttauth} & AugPAKE e \emph{tokens} por tópico
      & \Sim & \Nr & \Sim & \Sim & \Nr & \Nr & \Nr & \Nr & \Nr & \Nao & A, E & \Sim \\
    \multicolumn{15}{@{}l}{\textit{C4 --- Cifra baseada em atributos}} \\
    \E{01} & Singh et al.~\cite{singh2015secure}            & KP/CP-ABE sobre ECC
      & \Sim & \Nr & \Nr & \Sim & \Nr & \Nr & \Nr & \Nr & \Nr & \Nao & S & \Sim \\
    \multicolumn{15}{@{}l}{\textit{C5 --- Proteção leve fim-a-fim da mensagem}} \\
    \E{07} & Dinculeană e Cheng~\cite{dinculeana2019vulnerabilities} & Mapeamento \emph{Value-to-HMAC}
      & \Sim & \Sim & \Nr & \Nr & \Nr & \Nao & \Nr & \Nr & \Nr & \Nr & E & \Sim \\
    \E{08} & Malina et al.~\cite{malina2019secure}          & Três níveis de segurança (Schnorr)
      & \Parc & \Sim & \Sim & \Parc & \Parc & \Parc & \Sim & \Nr & \Sim & \Nao & A & \Sim \\
    \E{09} & Sadio et al.~\cite{sadio2019lightweight}       & AEAD ChaCha20-Poly1305
      & \Sim & \Sim & \Nr & \Nr & \Nr & \Nr & \Nr & \Nr & \Nr & \Sim & E & \Sim \\
    \E{13} & Dahlmanns et al.~\cite{dahlmanns2021entrust}   & Servidor de chaves por tópico (ENTRUST)
      & \Sim & \Sim & \Sim & \Sim & \Nr & \Sim & \Nr & \Parc & \Nr & \Sim & E & \Sim \\
    \E{16} & Van Glabbeek et al.~\cite{vanglabbeek2022flexible} & Três opções com assinatura Schnorr
      & \Parc & \Sim & \Sim & \Nr & \Parc & \Parc & \Sim & \Nr & \Sim & \Parc & E & \Sim \\
    \multicolumn{15}{@{}l}{\textit{C6 --- Proteção do broker e do transporte}} \\
    \E{12} & Segarra et al.~\cite{segarra2020mqttz}         & \emph{Broker} em TEE (TrustZone)
      & \Sim & \Nr & \Sim & \Parc & \Nr & \Nr & \Nr & \Nr & \Nr & \Parc & E & \Nr \\
    \E{14} & Spielvogel et al.~\cite{spielvogel2021tls}     & Nível de segurança exigido por salto
      & \Parc & \Nr & \Nr & \Nr & \Nr & \Nr & \Nr & \Nr & \Nr & \Nao & E & \Nr \\
    \E{15} & Prantl et al.~\cite{prantl2021performance}     & Avaliação do MQTT sobre TLS
      & \Sim & \Sim & \Nr & \Nr & \Nr & \Nr & \Nr & \Nr & \Nr & \Nao & E & \Sim \\
    \multicolumn{15}{@{}l}{\textit{C7 --- Deteção de intrusões}} \\
    \E{11} & Vaccari et al.~\cite{vaccari2020mqttset}       & Conjunto de dados para IDS
      & \Nr & \Nr & \Nr & \Nr & \Nr & \Nr & \Nr & \Parc & \Nr & \Nr & E & \Nr \\
    \E{18} & Alotaibi et al.~\cite{alotaibi2024secure}      & IDS distribuído (H2O, XGBoost)
      & \Nr & \Nr & \Nr & \Nr & \Nr & \Nr & \Nr & \Parc & \Nr & \Nr & E & \Nr \\
    \multicolumn{15}{@{}l}{\textit{C8 --- Privacidade e anonimato}} \\
    \E{17} & Buccafurri et al.~\cite{buccafurri2023mqtta}   & Rede de \emph{bridges} anónima
      & \Nr & \Nr & \Nr & \Nr & \Nr & \Nr & \Nr & \Nr & \Sim & \Nr & E & \Nr \\
    \bottomrule
  \end{tabular}
\end{table}

\paragraph{Cobertura das propriedades (QE3).} A Figura~\ref{fig:propriedades}
contabiliza, para os 16 mecanismos, quantos garantem cada propriedade (total ou
parcialmente). A confidencialidade é a propriedade mais tratada (12 estudos),
seguida da autorização (7), da autenticação (7) e da integridade (7). Só quatro
mecanismos funcionam, total ou parcialmente, com um \emph{broker} não confiável. A
validade temporal, a proteção contra repetição e o não repúdio são raros, mas não
estão ausentes: o ENTRUST deteta duplicados com um contador por
publicador~\cite{dahlmanns2021entrust}, e Malina et al.\ e Van Glabbeek et al.\
assinam um carimbo temporal, o que dá não repúdio e uma proteção parcial contra
repetição~\cite{malina2019secure, vanglabbeek2022flexible}. Nenhum dos três combina
um limite de idade com a deteção de duplicados dentro desse limite.

\begin{figure}[ht]
  \centering
  \begin{tikzpicture}
    \begin{axis}[
        xbar stacked, bar width=9pt, width=0.8\textwidth, height=6.4cm,
        xmin=0, xmax=18, xtick={0,2,...,18},
        symbolic y coords={NR,R,V,D,P,BnC,I,Az,Au,C},
        ytick=data, y tick label style={font=\small},
        xlabel={N.º de mecanismos (de 16)}, xlabel style={font=\small},
        legend style={font=\scriptsize, at={(0.98,0.04)}, anchor=south east},
        axis lines*=left, xmajorgrids, grid style={cinza!25}]
      \addplot[fill=acento!80, draw=iscte] coordinates {
        (7,C) (6,Au) (5,Az) (6,I) (2,BnC) (3,P) (0,D) (0,V) (1,R) (2,NR)};
      \addplot[fill=parcial!60, draw=parcial] coordinates {
        (5,C) (1,Au) (2,Az) (1,I) (2,BnC) (0,P) (3,D) (2,V) (2,R) (0,NR)};
      \legend{Garantida, Parcialmente}
    \end{axis}
  \end{tikzpicture}
  \caption{Número de mecanismos que garantem cada propriedade de segurança
    (siglas como na Tabela~\ref{tab:comparacao}).}
  \label{fig:propriedades}
\end{figure}

\paragraph{Avaliação (QF4).} Catorze dos 16 mecanismos ($87{,}5\,\%$) reportam
avaliação experimental; um é avaliado só em simulação~\cite{singh2015secure}, um só
analiticamente, por contagem de operações~\cite{malina2019secure}, e só um apresenta
uma prova formal de segurança~\cite{javadi2026cmap}. Onze são desenhados ou
avaliados para dispositivos restritos, do Arduino UNO~\cite{sadio2019lightweight} e
do Cortex-M3~\cite{dahlmanns2021entrust} ao Raspberry Pi~\cite{javadi2026cmap}. Nenhum dos estudos analisados mede, no mesmo ensaio, o custo de cada mecanismo e os
ataques que ele impede.

% ---------------------------------------------------------------------------
\subsection{Respostas às perguntas de investigação}
\label{sec:respostas}

\paragraph{QG e QF2.} Compreender as ameaças, proteger a ligação e o
\emph{broker}, e proteger a mensagem, sendo esta última a frente mais numerosa
(Secção~\ref{sec:taxonomia}).

\paragraph{QF1.} Escuta e injeção quando o MQTT é usado sem cifra, o que acontece
por omissão~\cite{andy2017attack}, negação de serviço que a cifra não
evita~\cite{vaccari2020slowite}, erros de implementação~\cite{hernandez2018fuzzing}
e o \emph{broker} como ponto de exposição de todo o
conteúdo~\cite{segarra2020mqttz, dahlmanns2021entrust}.

\paragraph{QF3 e QE3.} Ver a Figura~\ref{fig:propriedades}: a confidencialidade
domina, só quatro mecanismos dispensam a confiança no \emph{broker}, e a validade
temporal, a anti-repetição e o não repúdio aparecem isolados em três estudos.

\paragraph{QF4.} Sobretudo experimental (14 de 16); uma única prova formal.

\paragraph{QF5.} O custo da segurança em dispositivos a bateria~\cite{hintaw2023mqtt},
a inadequação do TLS a nós restritos~\cite{sadio2019lightweight, javadi2026cmap},
a exposição do conteúdo no \emph{broker}~\cite{segarra2020mqttz} e a impossibilidade
de um cliente impor a proteção dos restantes saltos~\cite{spielvogel2021tls}.
